\renewcommand{\thetheorem}{2.\arabic{theorem}}
\setcounter{theorem}{0}
\setcounter{chapter}{1}


\chapter{Logic and Mathematics}

Before advancing into the formal theory of vector spaces, we must establish the rigorous grammar of mathematical reasoning. This chapter provides the logical foundation necessary for constructing and understanding proofs.

\subsection{Mathematical Statements and Basic Operators}

\begin{definition}[Mathematical Statement]
A \newterm{mathematical statement} (or logical statement) is a declarative sentence that is definitively either \textbf{True (T)} or \textbf{False (F)}. It is a fundamental axiom of logic that a statement cannot be both true and false simultaneously.
\end{definition}

\begin{example}[Examples of Logical Statements]
Consider the following declarations:
\begin{enumerate}[label=\textbf{\Alph*}:]
    \item \qt{For every real number $x$, $x^2 \ge 0$.}
    \item \qt{Every prime number $p \ge 3$ is odd.}
    \item \qt{Every odd number $p \ge 3$ is prime.}
\end{enumerate}
Statements A and B are True. Statement C is False (e.g., $9$ is odd but not prime).
\end{example}

\subsection{Negation, Conjunction, and Disjunction}

\begin{definition}[Negation]
The \newterm{negation} of a statement $A$, denoted $\neg A$ (or sometimes $\sim A$), asserts that $A$ does not hold.
\[
\begin{array}{|c|c|}
\hline
A & \neg A \\
\hline
\text{T} & \text{F} \\
\text{F} & \text{T} \\
\hline
\end{array}
\]
\end{definition}

\begin{definition}[Conjunction / AND]
The statement $A \wedge B$ (read \qt{$A$ and $B$}) is True if and only if both $A$ and $B$ are True.
\[
\begin{array}{|c|c|c|}
\hline
A & B & A \wedge B \\
\hline
\text{T} & \text{T} & \text{T} \\
\text{T} & \text{F} & \text{F} \\
\text{F} & \text{T} & \text{F} \\
\text{F} & \text{F} & \text{F} \\
\hline
\end{array}
\]
\end{definition}

\begin{example}[Conjunction Example]
For $x \in \mathbb{R}$, the statement $(x^2 = 1) \wedge (x \ge 0)$ is logically equivalent to $x = 1$.
\end{example}

\begin{definition}[Disjunction / OR]
The statement $A \vee B$ (read \qt{$A$ or $B$}) is True if \textit{at least one} of the statements $A$ or $B$ is True. This is an \textit{inclusive} OR.
\[
\begin{array}{|c|c|c|}
\hline
A & B & A \vee B \\
\hline
\text{T} & \text{T} & \text{T} \\
\text{T} & \text{F} & \text{T} \\
\text{F} & \text{T} & \text{T} \\
\text{F} & \text{F} & \text{F} \\
\hline
\end{array}
\]
\end{definition}

\begin{example}[Disjunction Example]
The statement $(x > 1) \vee (x < 2)$ is True for all $x \in \mathbb{R}$. Conversely, $(x < 1) \vee (x > 2)$ is precisely the negation of the statement $1 \le x \le 2$.
\end{example}

\subsection{Implications and Equivalences}

\begin{definition}[Implication]
The statement $A \implies B$ (\qt{$A$ implies $B$}) is defined to be False only when $A$ is True and $B$ is False. In all other cases, it is True.
\[
\begin{array}{|c|c|c|}
\hline
A & B & A \implies B \\
\hline
\text{T} & \text{T} & \text{T} \\
\text{T} & \text{F} & \text{F} \\
\text{F} & \text{T} & \text{T} \\
\text{F} & \text{F} & \text{T} \\
\hline
\end{array}
\]
\end{definition}

\begin{remark}
Note that $A \implies B$ is logically equivalent to $\neg A \vee B$. Consequently, if the premise $A$ is False, the implication is \textit{vacuously true} regardless of $B$. As Prof. Biran noted:
\[
(0=1) \implies \qt{\text{The Earth is flat}}
\]
is a mathematically \textbf{True} statement. Causality is irrelevant in formal logic.
\end{remark}

\begin{definition}[Equivalence]
The statement $A \iff B$ (\qt{$A$ if and only if $B$}) holds when $A \implies B$ and $B \implies A$ are both True. This means $A$ and $B$ share the same truth value.
\[
\begin{array}{|c|c|c|}
\hline
A & B & A \iff B \\
\hline
\text{T} & \text{T} & \text{T} \\
\text{T} & \text{F} & \text{F} \\
\text{F} & \text{T} & \text{F} \\
\text{F} & \text{F} & \text{T} \\
\hline
\end{array}
\]
\end{definition}

\begin{exercise}[Contrapositive and De Morgan's Laws]
Prove the following via truth tables:
\begin{enumerate}
    \item \textbf{Contrapositive:} $(A \implies B) \iff (\neg B \implies \neg A)$.
    \item \textbf{De Morgan's Laws:} $\neg(A \wedge B) \iff \neg A \vee \neg B$ and $\neg(A \vee B) \iff \neg A \wedge \neg B$.
\end{enumerate}
\end{exercise}

\subsection{Predicates and Quantifiers}

A \newterm{predicate} $P(x)$ is a statement whose truth depends on a variable $x$. We use quantifiers to convert predicates into absolute statements.

\begin{definition}[Quantifiers]
\begin{itemize}
    \item \textbf{Universal ($\forall$):} \qt{For every $x$ in $X$, $P(x)$ holds.}
    \item \textbf{Existential ($\exists$):} \qt{There exists at least one $x$ in $X$ such that $P(x)$ holds.}
\end{itemize}
\end{definition}

\begin{example}[Quantifiers over Integers]
Let $P(n)$ be the predicate $n = n^2$ over $\mathbb{Z}$. 
\begin{itemize}
    \item $\forall n \in \mathbb{Z} : P(n)$ is False.
    \item $\exists n \in \mathbb{Z} : P(n)$ is True (take $n=1$).
\end{itemize}
\end{example}

\begin{importantremark}[The Order of Quantifiers]
The order of quantifiers is non-commutative. Consider $X$ (ETH students), $Y$ (courses), and $A(x,y) := \qt{\text{student } x \text{ takes course } y}$.
\begin{itemize}
    \item $\forall x \in X, \exists y \in Y : A(x,y)$ means every student takes at least one course (True).
    \item $\exists y \in Y, \forall x \in X : A(x,y)$ means there is one specific course taken by \textit{every} student (False).
\end{itemize}
\end{importantremark}

\subsection{Negation and Unique Existence}

\begin{proposition}[Negating Quantifiers]
\begin{align*}
\neg (\forall x \in X : A(x)) &\iff \exists x \in X : \neg A(x) \\
\neg (\exists x \in X : A(x)) &\iff \forall x \in X : \neg A(x)
\end{align*}
\end{proposition}

\begin{remark}[Shorthand Notation]
In practice, we often condense multiple quantifiers of the same type:
\begin{itemize}
    \item $\exists x \in X, \exists y \in X \implies \exists x,y \in X$
    \item $\forall x \in X, \forall y \in X \implies \forall x,y \in X$
\end{itemize}
\end{remark}

\begin{definition}[Unique Existence]
The notation $\exists! x \in X : P(x)$ means there exists exactly one $x$ satisfying the predicate. Formally:
\[
(\exists x \in X : P(x)) \wedge (\forall x, y \in X : (P(x) \wedge P(y) \implies x = y))
\]
\end{definition}

\begin{example}[Unique Existence of Roots]
$\exists! x \in \mathbb{R} : x^2 = 4$ is False (since $x = \pm 2$). However, $\exists! x \in \mathbb{R}_{\ge 0} : x^2 = 4$ is True.
\end{example}
